% ============================================================
% tex/preamble.tex —— 全局排版设置
%
% 只负责"整份文档长什么样"：页面、颜色、章节标题格式、
% 代码样式、页眉页脚、超链接。
% 正文命令在 tex/commands.tex，正文内容在 chapters/。
% ============================================================

% ---------- 中文数字章号：一、二、……十一 ----------
% 自己定义而不依赖 ctex 的 \chinese / zhnumber，避免宏包版本差异导致编译失败。
\newcommand{\chapterlabel}{%
  \ifcase\value{chapter}\or
  一\or 二\or 三\or 四\or 五\or 六\or
  七\or 八\or 九\or 十\or 十一\or 十二\fi
}

% ---------- 页面尺寸与边距 ----------
\usepackage{geometry}
\geometry{
  a4paper,
  top=2.8cm, bottom=2.5cm,
  left=3.0cm, right=3.0cm,
}

\linespread{1.35}

% ---------- 图片与表格 ----------
\usepackage{graphicx}
\usepackage{booktabs}
\usepackage{tabularx}
\usepackage{array}
\usepackage{caption}

% 长表格：xltabular = tabularx 的自动列宽 X + longtable 的跨页断行。
% 用于行数多到一页放不下的表（如第二章的场景清单），
% 否则浮动体被推迟到下一页，小标题会和表格分家。
\usepackage{xltabular}
\setlength{\LTpre}{8pt}    % 表前距（默认偏大）
\setlength{\LTpost}{8pt}   % 表后距
\setlength{\LTcapwidth}{\linewidth}
\renewcommand{\arraystretch}{1.15}
\captionsetup{font=small, labelsep=quad, skip=6pt}

% 浮动体放置参数。默认值下，同一节里挨着的两张表很容易被 LaTeX 判定为
% "这一页放不下正文了"，于是单独排成一页浮动页，正文被推到前后两页，
% 中间留出半页空白（第三章的色板表与字体表、视频表与转场表都撞上过）。
% 放宽"一页能放几个浮动体"和"浮动体最多占多大比例"，让表和正文同页。
\setcounter{topnumber}{3}
\setcounter{bottomnumber}{2}
\setcounter{totalnumber}{4}
\renewcommand{\topfraction}{0.92}      % 顶部浮动体最多占版心的比例
\renewcommand{\bottomfraction}{0.85}
\renewcommand{\textfraction}{0.06}      % 一页至少留多少正文
\renewcommand{\floatpagefraction}{0.85} % 不到这个比例就不单独排浮动页

% ---------- 流程图 ----------
% 用 TikZ 画在图环境里，不依赖外部工具导出的图片：改节点文字不用重新画图，
% 编译出来是矢量，缩放不糊。arrows.meta 提供 -{Stealth} 这类箭头样式。
\usepackage{tikz}
\usetikzlibrary{arrows.meta, positioning}

% ---------- 颜色 ----------
\usepackage{xcolor}
\definecolor{nightblue}{RGB}{26,35,58}    % 夜色深蓝：标题主色
\definecolor{lampgold}{RGB}{201,146,42}   % 灯焰暖黄：强调色
\definecolor{papergray}{RGB}{248,248,246} % 纸张灰：代码底色

% ---------- 代码块 ----------
\usepackage{listings}
\lstset{
  basicstyle=\small\ttfamily,
  backgroundcolor=\color{papergray},
  commentstyle=\color{black!45},
  keywordstyle=\color{nightblue}\bfseries,
  stringstyle=\color{lampgold},
  numbers=left,
  numberstyle=\tiny\color{black!35},
  numbersep=8pt,
  frame=single,
  framerule=0.4pt,
  rulecolor=\color{black!15},
  breaklines=true,
  showstringspaces=false,
  tabsize=2,
  columns=flexible,
  keepspaces=true,
}

% ---------- 章节标题格式 ----------
\ctexset{
  chapter = {
    name       = {},
    number     = \chapterlabel,
    aftername  = {、},
    format     = \centering\Large\bfseries\color{nightblue},
    beforeskip = 1.0em,
    afterskip  = 1.8em,
    pagestyle  = fancy,
  },
  section = {
    format     = \large\bfseries\color{nightblue},
    aftername  = {\quad},
    beforeskip = 1.2em,
    afterskip  = 0.8em,
  },
  subsection = {
    format = \normalsize\bfseries\color{black!85},
  },
  subsubsection = {
    format = \normalsize\bfseries\color{black!75},
  },
}

% ---------- 目录里章号的分隔符 ----------
\usepackage{etoolbox}

% ctex 写入 .toc 时把章号与标题之间的分隔符硬编码成 \hspace{.3em}，
% 与标题行的「、」对不上，这里改掉。补丁失败会打印警告，不会被无声忽略。
\makeatletter
\AtBeginDocument{%
  \patchcmd{\CTEX@chapter@tocline}
    {\hspace{.3em}}
    {、}
    {}{\PackageWarning{project-report}{目录章号分隔符补丁未生效：TOC 里会显示成“一 标题”}}%
}
\makeatother

% ---------- 页眉页脚 ----------
\usepackage{fancyhdr}
\setlength{\headheight}{15pt}
\pagestyle{fancy}
\fancyhf{}
\fancyhead[L]{\small\color{black!50}\ProjectName · 项目文档}
\fancyhead[R]{\small\color{black!50}\leftmark}
\fancyfoot[C]{\small\thepage}
\renewcommand{\headrulewidth}{0.4pt}
\renewcommand{\footrulewidth}{0pt}

% 章首页沿用相同页眉页脚（否则章首页会退回朴素样式）
\fancypagestyle{plain}{
  \fancyhf{}
  \fancyhead[L]{\small\color{black!50}\ProjectName · 项目文档}
  \fancyhead[R]{\small\color{black!50}\leftmark}
  \fancyfoot[C]{\small\thepage}
  \renewcommand{\headrulewidth}{0.4pt}
  \renewcommand{\footrulewidth}{0pt}
}

% ---------- 超链接（必须最后加载）----------
\usepackage{hyperref}
\hypersetup{
  colorlinks        = true,
  linkcolor         = nightblue,
  citecolor         = lampgold,
  urlcolor          = lampgold,
  bookmarksnumbered = true,
  bookmarksopen     = true,
  pdftitle          = {\ProjectName 项目文档},
  pdfauthor         = {\TeamName},
}
